% English version of paper/ja_v4/05_考察.md (ver4, 2026-11-17). MIT CommLab discussion: main conclusion -> answers to the questions of the introduction -> agreement and disagreement with prior work -> what the limitations leave open -> next steps. Cited claims as in reviews/v4_ja/citations_2026-11-17.md.
\FloatBarrier
\section{Discussion}\label{sec:discussion}

The main conclusion of this study is that, on a new corpus with few labels, the other modalities of the same utterance made up for part of the missing labels, and how much they made up depended on how the pseudo-labels were built. Albanie et al.\ \cite{albanie2018} and Zhang et al.\ \cite{zhang2021combining} showed that classifiers can be trained on predictions of other modalities; we showed, by comparison under the same conditions in a cross-corpus setting with only a few target labels, that doing so gave audio classifiers higher UAR than domain adaptation, semi-supervised and test-time adaptation methods.

The answers to the two questions of the introduction are as follows. First, on whether same-modality methods using unlabelled utterances or pseudo-labels including other modalities make better classifiers, pseudo-labels including other modalities gave audio classifiers higher UAR (Sect.~\ref{sec:res-methods}). In the all-source setting, same-modality unlabelled utterances rarely reduced the remaining error, whereas pseudo-labels including other modalities reduced it in nearly half of the conditions (Sect.~\ref{sec:res-remaining}). One possible explanation is that same-modality unlabelled utterances add only information on the distribution of features the classifier already uses, whereas predictions of other modalities add separate cues about the emotion of the utterance. Another possible explanation of why distillation gave audio classifiers gains with intervals above zero in 9 of 12 conditions is that audio features predicted the decisions of the text classifier, which is part of the product, better than visual features did (Supplementary Information S3; its relation to the gains was not measured). Second, on how to build the pseudo-labels, retrained classifiers had higher UAR the more accurate the pseudo-labels were relative to the classifier before retraining (Sect.~\ref{sec:res-margin}), and setting per-emotion weights of each modality from the few labels and the sources produced pseudo-labels more accurate than equal weights, with higher UAR also for classifiers retrained on them (Sects.~\ref{sec:res-prop}--\ref{sec:res-retrain}).

Our results agree and disagree with prior work in the following points. That distillation gave audio classifiers higher UAR is in the same direction as Albanie et al.\ \cite{albanie2018}, who showed that audio emotion classifiers can be trained on facial predictions; we confirmed it in a cross-corpus setting with only a few target labels. Training on the predictions of a more accurate model of another modality has been reported not to always improve the trained model \cite{xue2023mfh}. In our results, no variant with a negative advantage, that is, with pseudo-labels less accurate than the classifier before retraining, rose, and UAR rose for only some variants with more accurate pseudo-labels (Sect.~\ref{sec:res-margin}). This suggests that pseudo-labels being no less accurate than the classifier before retraining is a necessary, but not a sufficient, condition for UAR to rise. The emotion-dependent differences in modality strength that Busso et al.\ \cite{busso2004} reported within one corpus were also observed across our corpora, where the ranking of classifier strengths changed from corpus to corpus (Sect.~\ref{sec:idea}). Saito et al.\ \cite{saito2019mme} reported that, when the target has a few labels, unsupervised domain adaptation methods that align feature distributions often do not improve performance relative to just training on labelled source and target examples. In our all-source setting, the eight methods, including versions of DANN and MME moved to linear classifiers on frozen features, rarely reduced the remaining error (Sect.~\ref{sec:res-remaining}). With few training samples, a model trained on more modalities has also been argued to be able to fall below a model trained on a subset of them \cite{huang2021}. The proposal combines the decisions of three separately trained classifiers with three quantities set on the K labels, and it had higher UAR than the single classifier strongest on the K labels in all 12 conditions (Sect.~\ref{sec:res-prop}).

Several points remain open within the scope of this study (Sect.~\ref{sec:limitations}). The classifiers are linear classifiers on frozen encoders, and whether the comparisons give the same results with fine-tuned encoders is unknown. Distillation gave audio classifiers gains with intervals above zero in 9 of 12 conditions, whereas the gains of the visual, text and AVT classifiers varied across corpora. The proposal was built on the same four corpora used to evaluate it, so differences from the product and other combinations may be biased in its favour. The relation between the accuracy of the pseudo-labels and UAR is a descriptive analysis set up after the computation.

We plan three next steps. First, we will compare same-modality methods with pseudo-labels from other modalities using fine-tuned encoders and other corpora. Second, we will evaluate the proposal once on a new corpus with the procedure fixed and retrain the audio, visual and text classifiers on its pseudo-labels. Third, we will estimate from the few labels whether the pseudo-labels are more accurate than the classifier before retraining, and turn this into a procedure for choosing pseudo-label providers. Building MER classifiers for a new corpus involves weighing the effort of labelling against that of collecting unlabelled multimodal recordings; as material for this decision, our results showed how much pseudo-labels including other modalities made up for the shortage of labels and that this depended on how the pseudo-labels were built.
